%==============================================================================
%  MATH 231 -- Linear Algebra I (Queens College, Fall 2026)
%  PLAN / NOT FINALIZED -- Computational projects, application demos, and a
%  draft assessment mapping.
%
%  Split out of MATH231_progression so the progression document states what the
%  course covers rather than what it might.  Which projects get assigned is
%  still open.
%
%  NOTE: the assessment mapping below predates the finalized grading split and
%  disagrees with it (4 projects / 10 homework here vs. 3 projects / 5 homework
%  actual).  Kept for the unit-coverage sketch only.  See plan/README.md.
%==============================================================================
\documentclass[11pt]{article}

\usepackage[margin=1in]{geometry}
\usepackage{amsmath,amssymb}
\usepackage[dvipsnames]{xcolor}
\usepackage{enumitem}
\usepackage{booktabs}
\usepackage{array}
\usepackage[hidelinks]{hyperref}
\usepackage{titlesec}

\titleformat{\section}{\Large\bfseries\sffamily\color{RoyalBlue}}{}{0pt}{}
\titleformat{\subsection}{\large\bfseries\sffamily}{}{0pt}{}
\titlespacing*{\section}{0pt}{18pt}{8pt}

\newcommand{\pri}[1]{\textcolor{BrickRed}{\small\sffamily[#1]}}

\title{\vspace{-1.2cm}\textbf{\sffamily MATH 231 -- Linear Algebra I}\\[4pt]
       {\large\sffamily Computational Projects \& Application Demos}\\[2pt]
       {\normalsize\sffamily Fall 2026 \textbullet\ working draft --- not finalized}}
\author{}
\date{}

\begin{document}
\maketitle
\vspace{-1.3cm}

\begin{center}
\fcolorbox{BrickRed!50}{BrickRed!7}{%
\begin{minipage}{0.93\textwidth}
\small\sffamily
\textbf{This document is not finalized.} It is a menu of candidates, not a list
of assignments. Nothing here is part of the syllabus, and no item should be
treated as assigned unless it has been officially announced. Where this document
and an official course document disagree, the official document is correct.
\end{minipage}}
\end{center}

\bigskip

%==============================================================================
\section{Assessment mapping (draft)}
%==============================================================================
\noindent\emph{Counts here are stale:} this sketch shows four projects and ten
homework assignments, while the course carries \textbf{three projects and five
homework assignments}. Retained for the unit-coverage mapping only; the root
\texttt{README.md} holds the authoritative weights.

\begin{center}
\renewcommand{\arraystretch}{1.3}
\begin{tabular}{@{}l p{4.0cm} p{6.4cm}@{}}
\toprule
\textbf{Assessment} & \textbf{Units covered} & \textbf{Suggested focus} \\
\midrule
Exam 1              & 1--6    & vector spaces, norms, complex, matrices, calculus, conditioning \\
Exam 2              & 7--10   & systems, maps, subspaces, least squares \\
Final Exam & 1--12 (weighted to 11--12) & eigen/SVD, optimization, synthesis \\
\midrule
Project 1 & after 2  & \#1 K-means image color segmentation \\
Project 2 & after 11 & \#5 Eigencats via PCA \\
Project 3 & after 12 & \#7 Mini-ML: OLS + SVM + KNN with validation \\
Project 4 & after 12 & \#8 Signal trilateration: Gauss--Newton / LM \\
Lecture demos & 10--12 & \#2 SVD+DCT, \#3 LSA, \#4 collab.\ filtering, \#6 PageRank, \#9 embeddings \\
\midrule
Homework & one per unit (Units 1--12, some combined) & mixed hand computation + Python verification \\
\bottomrule
\end{tabular}
\end{center}

%==============================================================================
\section{The project menu}
%==============================================================================
This is a \emph{menu}, not an assignment quota: every item below is documented,
but only a small subset is assigned to students as an implementation; the
complement are instructor-led demos that still put the application in front of
the class. The list is open --- further items may be appended as \#10, \#11,
\dots\ without disturbing the ordering.

\smallskip
\noindent\emph{Outline additions these items introduced:} DCT / orthogonal
transforms (Unit~10); Perron--Frobenius theorem and stochastic matrices
(Unit~11); TF--IDF term--document weighting (Unit~11); masked-loss / ALS
factorization and train/validation/test model selection (Unit~12).

\begin{center}
\renewcommand{\arraystretch}{1.25}
\begin{tabular}{@{}c l c c p{4.6cm}@{}}
\toprule
\textbf{\#} & \textbf{Project} & \textbf{Role} & \textbf{Prereq} & \textbf{LA machinery} \\
\midrule
1 & K-means color segmentation   & Assigned & U2  & K-means, image vectorization, $\ell_2$ distance \\
2 & SVD + DCT image compression  & Demo     & U11 & SVD, Eckart--Young, DCT orthonormal basis \\
3 & LSA file-system search       & Demo     & U11 & truncated SVD, TF--IDF, cosine similarity \\
4 & Collaborative filtering      & Demo     & U12 & masked low-rank factorization, ALS/SGD \\
5 & Eigencats (PCA)              & Assigned & U11 & PCA, covariance eigenbasis / SVD \\
6 & PageRank                     & Demo     & U11 & power method, stochastic matrix, Perron--Frobenius \\
7 & Mini-ML (OLS/SVM/KNN)        & Assigned & U12 & OLS, SVM, KNN, train/val/test \\
8 & Signal trilateration         & Assigned & U12 & Gauss--Newton, Levenberg--Marquardt, Jacobian \\
9 & Semantic embedding geometry  & Demo     & U11 & cosine similarity, KNN, PCA, vector arithmetic \\
\bottomrule
\end{tabular}
\end{center}

\noindent All implementations are in \textbf{Python}.

%==============================================================================
\section{Candidate assigned implementations}
%==============================================================================
\begin{description}[leftmargin=0pt,style=nextline,font=\normalfont\sffamily\bfseries\color{RoyalBlue},itemsep=6pt]
\item[\#1 -- K-means image color segmentation (after Unit 2)]
  Cluster an image's pixels in RGB (or CIELAB) space into $k$ colors and replace
  each pixel by its cluster centroid; display the quantized/segmented image as a
  function of $k$. \emph{LA focus:} $\ell_2$ distance, centroids as means,
  Lloyd's iteration and its convergence. Python (NumPy/Pillow).
\item[\#5 -- Eigencats via PCA (after Unit 11)]
  Assemble a data matrix of aligned cat (or face) images, mean-center, and
  compute principal components via the SVD; visualize the top ``eigencats,''
  reconstruct images from $r$ components, and plot reconstruction error and
  variance-explained vs.\ $r$. \emph{LA focus:} covariance eigenstructure, SVD,
  low-rank reconstruction. Python.
\item[\#7 -- Mini-ML with validation (after Unit 12)]
  Three sub-tasks on suitable datasets: OLS for a regression problem, an SVM for
  a classification problem, and KNN for a multiclass problem; use a
  train/validation/test split, report generalization metrics, and justify model
  choice. \emph{LA focus:} normal equations vs.\ QR (OLS), margin/kernels (SVM),
  distance metrics (KNN), validation methodology. Python (hand-implement the
  cores; library baselines allowed).
  \emph{Depends on topics currently listed as not guaranteed --- see
  \texttt{MATH231\_tentative\_topics}.}
\item[\#8 -- Signal trilateration (after Unit 12)]
  From noisy range measurements to known anchors, estimate an emitter's position
  by nonlinear least squares; implement Gauss--Newton and Levenberg--Marquardt
  and compare convergence/robustness vs.\ initialization and noise.
  \emph{LA focus:} residual Jacobian, GN normal-equation step, LM damping,
  conditioning. Python.
\end{description}

%==============================================================================
\section{Candidate lecture demos}
%==============================================================================
\begin{description}[leftmargin=0pt,style=nextline,font=\normalfont\sffamily\bfseries\color{RoyalBlue},itemsep=6pt]
\item[\#2 -- SVD + DCT image compression (Unit 11)]
  Sweep the rank $r$ of an image's SVD and watch quality vs.\ storage; then show
  the $8\times8$ block DCT that JPEG actually uses, connecting a \emph{fixed}
  orthonormal basis to the data-adaptive SVD basis. Motivates Eckart--Young.
\item[\#3 -- Latent Semantic Analysis search (Unit 11)]
  Build a small term--document matrix over a mock file system, TF--IDF weight it,
  truncate its SVD, and answer queries by cosine similarity in the reduced
  ``concept'' space; contrast with exact keyword matching.
\item[\#4 -- Collaborative filtering (Unit 12)]
  On a toy ratings matrix with missing entries, fit a low-rank factorization by
  ALS/SGD over \emph{observed} entries and predict held-out ratings; frames
  recommendation as matrix completion.
\item[\#6 -- PageRank (Unit 11)]
  Model a small web graph as a stochastic matrix, run the power method, and watch
  the ranking vector converge to the dominant eigenvector; connect to
  Perron--Frobenius and the damping factor.
\item[\#9 -- Semantic embedding geometry (Unit 11, static)]
  With pretrained word vectors: nearest neighbors under cosine similarity,
  analogy arithmetic ($\text{king}-\text{man}+\text{woman}\approx\text{queen}$),
  PCA to $2$-D for visualization, and clustering --- purely linear algebra, no
  model training.
\end{description}

\end{document}
